\section{Розділ~\ref{sec:debugging}. Налагодження машини}

Що чує електрод, маловимірність активності, роботу інтерфейсів на жорстких ґратках, спільне навчання мозку й декодера та зорові точки від стимуляції виміряно в рецензованих дослідах; оцінки потоків даних --- арифметика розділу; про вживлений людям пристрій Neuralink, наскільки вдалося перевірити, публікувала дані переважно сама компанія.

{\footnotesize
\setlength{\tabcolsep}{3pt}\renewcommand{\arraystretch}{1.15}
\begin{longtable}{>{\raggedright}p{0.5cm}>{\raggedright}p{4.0cm}>{\raggedright}p{5.6cm}>{\raggedright}p{2.3cm}>{\raggedright\arraybackslash}p{1.7cm}}
\toprule
№ & Твердження & Дослід і що виміряно & Джерело & Статус \\
\midrule\endfirsthead
\toprule
№ & Твердження & Дослід і що виміряно & Джерело & Статус \\
\midrule\endhead
1 & Електрод чує імпульси нейронів у радіусі порядку сотні мікрометрів, зазвичай від одного до кількох; їх розрізняють за формою & Щур, поле CA1, одночасний внутрішньоклітинний і позаклітинний запис: форма позаклітинного імпульсу повторює ширину й амплітуду внутрішньоклітинного. Оцінка: тетрод міг би розрізняти $60$--$100$ нейронів, а насправді виділяють близько шести. & \cite{Henze2000extra} & виміряно \\
2 & У мозку $8.6\cdot10^{10}$ нейронів; щуп чує приблизно одну стомільйонну & Підрахунок ядер у суспензії розчиненої тканини: $86\cdot10^9$ нейронів у мозку людини. Частка --- арифметика розділу. & \cite{Azevedo2009} & виміряно \\
3 & Активність сотень нейронів рухової кори вкладається в один-два десятки спільних змінних & Огляд записів у руховій корі мавп: активність популяції описується небагатьма прихованими змінними, спільними для багатьох нейронів. & \cite{Gallego2017} & виміряно \\
4 & Шум сусідів узгоджений, і сотня нейронів несе майже стільки ж, скільки тисяча (твердження~\ref{prop:pooling}) & Зорова кора мавпи: коефіцієнт кореляції шуму сусідніх нейронів близько $0.1$, і виграш від усереднення насичується на сотні нейронів. Теорія кореляцій, що обмежують інформацію. & \cite{Zohary1994, MorenoBote2014} & виміряно \\
5 & Сирий потік $2\cdot10^8$~біт/с проти $8\cdot10^4$~біт/с в імпульсах~\eqref{eq:probedata} & Арифметика розділу за ємністю потоку імпульсів~\eqref{eq:spikeentropy}. Параметри оцифрування реальні: кристал Neuralink --- $19.3$~кГц, $10$~біт на канал. & виведення в розділі; \cite{Rieke1997, Musk2019} & теорія \\
6 & Перша система Neuralink: кристали підсилюють і оцифровують, сирий потік іде кабелем, імпульси виділяє програма ззовні; виділення на кристалі, закритий корпус, радіо й бездротове живлення --- у пристрою для людей, за повідомленнями компанії & Стаття компанії: повний сирий потік кабелем USB-C, імпульси виділяє програма в реальному часі поза кристалом; стиснення на борту, бездротові живлення й передавання названо планом для клінічних пристроїв. Для вживленого людям пристрою --- лише повідомлення компанії. Що виділяти імпульси на кристалі необхідно, --- висновок розділу з~\eqref{eq:probedata}. & \cite{Musk2019}; звіти Neuralink, рецензованої статті немає & виміряно (стаття компанії); для людей --- звіт компанії \\
7 & До $3072$ електродів на $96$ нитках по $32$; нитки вводить робот, оминаючи судини & Стаття компанії: $3072$ електроди на $96$ нитках по $32$; робот уводить $6$ ниток ($192$ електроди) за хвилину; до $70\,\%$ електродів у щурів із довготривало вживленою ґраткою чують імпульси. & \cite{Musk2019} & виміряно (стаття компанії) \\
8 & У людей близько тисячі каналів на $64$ нитках; частина ниток змістилася; курсор порядку $10$~біт/с & Лише повідомлення компанії. Пошук у PubMed не знайшов рецензованої статті з результатами на людях; це узгоджується із застереженням у тексті розділу. & звіти Neuralink; рецензованої статті немає & виміряно (звіт компанії) \\
9 & Інтерфейс на ґратці із сотні електродів керує курсором & Людина з паралічем, $96$ електродів у первинній руховій корі: намір рухати рукою змінює імпульси через три роки після травми; декодер дає «нейронний курсор» (пошта, телевізор), керування протезом кисті й роботом-рукою. & \cite{Hochberg2006} & виміряно \\
10 & Письмо «уявною рукою»: близько $90$ знаків на хвилину, менше ніж $8$~біт/с & Людина з паралічем кисті, спроби писати від руки, рекурентний декодер: $90$ знаків на хвилину з точністю $94.1\,\%$ у реальному часі, понад $99\,\%$ після автокорекції. Оцінка в бітах --- арифметика розділу. & \cite{Willett2021} & виміряно \\
11 & Спроби говорити перетворюються на текст зі швидкістю близько $60$ слів на хвилину & Людина, що втратила розбірливе мовлення, ґратки в корі: $62$ слова на хвилину; $9.1\,\%$ помилок у словах за словника з $50$ слів, $23.8\,\%$ за словника зі $125\,000$. & \cite{Willett2023} & виміряно \\
12 & Інтерфейс витягає малу частку ємності: в одного нейрона десятки біт на секунду, у сотні --- тисячі (твердження~\ref{prop:probe}) & Верхня межа~\eqref{eq:spikeentropy} --- теорія; порівняння з рядками~8, 10 --- висновок розділу. Що вузьке місце --- маловимірність і незнання коду, а не провід, прямим дослідом не перевірено. & \cite{Rieke1997} & теорія \\
13 & Мозок і декодер учаться один в одного & Мавпи керують курсором; декодер підлаштовується в замкненій петлі, а нейрони водночас перенавчаються: точність зростає, навичка зберігається й менше страждає від рухів власної руки. Порада розвести швидкості навчання --- інженерне правило, окремого досліду в розділі немає. & \cite{Orsborn2014coad} & виміряно \\
14 & Струм через електрод збуджує навколишні нейрони й проводи без розбору типу & Кора миші й кішки, двофотонний запис кальцію: навіть слабкий струм збуджує рідкісні нейрони, розкидані на відстань до міліметрів, через пряме збудження аксонів в об'ємі діаметром у десятки мікрометрів. Тобто збудження невибіркове, але не суцільне. & \cite{Histed2009stim} & виміряно \\
15 & Стимуляція зорової кори запалює точки; до $16$ точок одночасно дають змогу впізнати деякі літери й межі предметів & Незряча людина, $96$ електродів у зоровій корі на $6$ місяців: поріг одного електрода $66.8\pm36.5$~мкА; одночасна стимуляція груп (до $16$ електродів --- межа стимулятора) знижує поріг і дає розрізненні візерунки, деякі літери й межі предметів упізнаються. & \cite{Fernandez2021} & виміряно \\
16 & Другий напрям Neuralink --- пристрій, що подає сигнали в зорову кору & Лише повідомлення компанії про розроблення; рецензованих даних про його роботу не знайдено. & звіти Neuralink & гіпотеза \\
17 & Концентрації широкомовних медіаторів можна вимірювати хімічним датчиком у тканині & Зрізи мозку щура, швидка циклічна вольтамперометрія на вуглецевому волокні: виміряно викид і зворотне захоплення серотоніну; речовина виходить за межі синапсу. & \cite{Bunin1998} & виміряно \\
\bottomrule
\end{longtable}}
